% GENERATED FILE. Edit canonical lessons, then run npm run generate:books.
% canonical-source-hash: fnv1a64:de6c51228be74610
% canonical-lessons: PA-C317-vela, PA-C317-bitna, PA-C317-zamana, PA-C317-pachhtauna, PA-C317-guzarna

\chapter{Days Gone By}
\label{ch:pa-a2-days-gone-by}
\hlchaptermodality{317}

\begin{chapteropening}
\textbf{By the end of this chapter:} \emph{I can say how time went by and how I felt about it.}

Complete the last lesson of chapter 317: I can say how time went by and how I felt about it.
\end{chapteropening}

\section[\texorpdfstring{\pa{ਵੇਲਾ} (velā)}{ਵੇਲਾ (velā)}]{\pa{ਵੇਲਾ} (velā) — a particular time}

\label{lesson:PA-C317-vela}



\begin{quote}
Before the new one: say the Punjabi for a swear word, then the Punjabi for a reconciliation.
\end{quote}

\subsection*{You'll want to know: \pa{ਵੇਲਾ}}
\textbf{\pa{ਵੇਲਾ}} — \emph{velā} — \textquotedblleft{}a particular time\textquotedblright{}.

\begin{quote}
\pa{ਉਸ} \pa{ਵੇਲੇ} \pa{ਮੈਂ} \pa{ਸਕੂਲ} \pa{ਵਿੱਚ} \pa{ਸੀ।}
\end{quote}

\begin{quote}
\emph{us vele maiṁ sakūl vich sī.}
\end{quote}

\begin{quote}
At that time I was at school.
\end{quote}

A masculine noun.

\begin{cousinweb}
\textbf{\pa{ਵੇਲਾ}} (\emph{velā}) is inherited from Sanskrit \emph{velā}, \textquotedblleft{}a limit, a time\textquotedblright{}; CDIAL lists Punjabi \emph{velā}, \textquotedblleft{}time\textquotedblright{}, under it (CDIAL 12115).
\end{cousinweb}

\begin{grammarlens}[title={What you've built}]
One more word for saying what happened.
\end{grammarlens}

\subsection*{Your turn}
\begin{itemize}
\raggedright
  \item \emph{Say it:} \emph{velā}
  \item \emph{Say it:} \emph{velā}, once more
  \item \emph{Say it:} \emph{us vele maiṁ sakūl vich sī.} — At that time I was at school.
  \item \emph{Recall:} say the Punjabi for a swear word, then the Punjabi for a reconciliation, then say \emph{velā} again
\end{itemize}

\begin{tcolorbox}[breakable,colback=teal!4,colframe=teal!35!black,title={Before you move on}]
What does \pa{ਵੇਲਾ} mean? (\textquotedblleft{}A particular time\textquotedblright{}.) Say it once more.
\end{tcolorbox}
\section[\texorpdfstring{\pa{ਬੀਤਣਾ} (bītṇā)}{ਬੀਤਣਾ (bītṇā)}]{\pa{ਬੀਤਣਾ} (bītṇā) — to go by (of time)}

\label{lesson:PA-C317-bitna}



\begin{quote}
Before the new one: say the Punjabi for a reconciliation, then the Punjabi for a particular time.
\end{quote}

\subsection*{You'll want to know: \pa{ਬੀਤਣਾ}}
\textbf{\pa{ਬੀਤਣਾ}} — \emph{bītṇā} — \textquotedblleft{}to go by (of time)\textquotedblright{}.

\begin{quote}
\pa{ਛੁੱਟੀ} \pa{ਦੇ} \pa{ਦਿਨ} \pa{ਬਹੁਤ} \pa{ਜਲਦੀ} \pa{ਬੀਤੇ।}
\end{quote}

\begin{quote}
\emph{chhuṭṭī de din bahut jaldī bīte.}
\end{quote}

\begin{quote}
The days of the holiday went by very quickly.
\end{quote}

Past (it …): \pa{ਉਹ} \pa{ਸਾਲ} \pa{ਬਹੁਤ} \pa{ਔਖਾ} \pa{ਬੀਤਿਆ}, \emph{oh sāl bahut aukhā bītiā}.

\begin{cousinweb}
\textbf{\pa{ਬੀਤਣਾ}} (\emph{bītṇā}) is inherited from Sanskrit \emph{vṛtta}, \textquotedblleft{}turned; passed (of time)\textquotedblright{}; CDIAL lists Punjabi \emph{batāuṇā}, \textquotedblleft{}to pass (time)\textquotedblright{}, under it, and gives the intransitive \emph{bītnā}, \textquotedblleft{}to pass (of time)\textquotedblright{}, only for Hindi (CDIAL 12069).
\end{cousinweb}

\begin{grammarlens}[title={What you've built}]
One more word for saying what happened.
\end{grammarlens}

\subsection*{Your turn}
\begin{itemize}
\raggedright
  \item \emph{Say it:} \emph{bītṇā}
  \item \emph{Say it:} \emph{bītṇā}, once more
  \item \emph{Say it:} \emph{chhuṭṭī de din bahut jaldī bīte.} — The days of the holiday went by very quickly.
  \item \emph{Recall:} say the Punjabi for a reconciliation, then the Punjabi for a particular time, then say \emph{bītṇā} again
  \item \emph{Write it:} copy \textbf{\pa{ਬੀਤਣਾ}} once, checking the whole word
\end{itemize}

\begin{tcolorbox}[breakable,colback=teal!4,colframe=teal!35!black,title={Before you move on}]
What does \pa{ਬੀਤਣਾ} mean? (\textquotedblleft{}To go by (of time)\textquotedblright{}.) Say it once more.
\end{tcolorbox}
\section[\texorpdfstring{\pa{ਜ਼ਮਾਨਾ} (zamānā)}{ਜ਼ਮਾਨਾ (zamānā)}]{\pa{ਜ਼ਮਾਨਾ} (zamānā) — an era}

\label{lesson:PA-C317-zamana}



\begin{quote}
Before the new one: say the Punjabi for a particular time, then the Punjabi for to go by (of time).
\end{quote}

\subsection*{You'll want to know: \pa{ਜ਼ਮਾਨਾ}}
\textbf{\pa{ਜ਼ਮਾਨਾ}} — \emph{zamānā} — \textquotedblleft{}an era\textquotedblright{}.

\begin{quote}
\pa{ਉਸ} \pa{ਜ਼ਮਾਨੇ} \pa{ਵਿੱਚ} \pa{ਲੋਕ} \pa{ਪੈਦਲ} \pa{ਜਾਂਦੇ} \pa{ਸਨ।}
\end{quote}

\begin{quote}
\emph{us zamāne vich lok paidal jānde san.}
\end{quote}

\begin{quote}
In those days people went on foot.
\end{quote}

A masculine noun.

\begin{cousinweb}
\textbf{\pa{ਜ਼ਮਾਨਾ}} (\emph{zamānā}) is a loan, usually traced through Persian to Arabic \emph{zamān}, \textquotedblleft{}time, a period\textquotedblright{}.
\end{cousinweb}

\begin{grammarlens}[title={What you've built}]
One more word for saying what happened.
\end{grammarlens}

\subsection*{Your turn}
\begin{itemize}
\raggedright
  \item \emph{Say it:} \emph{zamānā}
  \item \emph{Say it:} \emph{zamānā}, once more
  \item \emph{Say it:} \emph{us zamāne vich lok paidal jānde san.} — In those days people went on foot.
  \item \emph{Recall:} say the Punjabi for a particular time, then the Punjabi for to go by (of time), then say \emph{zamānā} again
  \item \emph{Write it:} copy \textbf{\pa{ਜ਼ਮਾਨਾ}} once, checking the whole word
\end{itemize}

\begin{tcolorbox}[breakable,colback=teal!4,colframe=teal!35!black,title={Before you move on}]
What does \pa{ਜ਼ਮਾਨਾ} mean? (\textquotedblleft{}An era\textquotedblright{}.) Say it once more.
\end{tcolorbox}
\section[\texorpdfstring{\pa{ਪਛਤਾਉਣਾ} (pachhtāuṇā)}{ਪਛਤਾਉਣਾ (pachhtāuṇā)}]{\pa{ਪਛਤਾਉਣਾ} (pachhtāuṇā) — to regret}

\label{lesson:PA-C317-pachhtauna}



\begin{quote}
Before the new one: say the Punjabi for to go by (of time), then the Punjabi for an era.
\end{quote}

\subsection*{You'll want to know: \pa{ਪਛਤਾਉਣਾ}}
\textbf{\pa{ਪਛਤਾਉਣਾ}} — \emph{pachhtāuṇā} — \textquotedblleft{}to regret\textquotedblright{}.

\begin{quote}
\pa{ਝਗੜੇ} \pa{ਤੋਂ} \pa{ਬਾਅਦ} \pa{ਮੈਂ} \pa{ਬਹੁਤ} \pa{ਪਛਤਾਇਆ।}
\end{quote}

\begin{quote}
\emph{jhagṛe ton bāad maiṁ bahut pachhtāiā.}
\end{quote}

\begin{quote}
After the quarrel I regretted it a lot.
\end{quote}

Past (I …): \pa{ਮੈਂ} \pa{ਬਾਅਦ} \pa{ਵਿੱਚ} \pa{ਪਛਤਾਇਆ}, \emph{maiṁ bāad vich pachhtāiā}.

\begin{cousinweb}
\textbf{\pa{ਪਛਤਾਉਣਾ}} (\emph{pachhtāuṇā}) goes back to a reconstructed \emph{*paścottāpa}, \textquotedblleft{}repentance\textquotedblright{}; CDIAL lists Punjabi \emph{pachotāvā}, \textquotedblleft{}repentance\textquotedblright{}, under it, beside the -ṇā forms \emph{pachõtāṇā} and \emph{pastāṇā} (CDIAL 8010).
\end{cousinweb}

\begin{grammarlens}[title={What you've built}]
One more word for saying what happened.
\end{grammarlens}

\subsection*{Your turn}
\begin{itemize}
\raggedright
  \item \emph{Say it:} \emph{pachhtāuṇā}
  \item \emph{Say it:} \emph{pachhtāuṇā}, once more
  \item \emph{Say it:} \emph{jhagṛe ton bāad maiṁ bahut pachhtāiā.} — After the quarrel I regretted it a lot.
  \item \emph{Recall:} say the Punjabi for to go by (of time), then the Punjabi for an era, then say \emph{pachhtāuṇā} again
  \item \emph{Write it:} copy \textbf{\pa{ਪਛਤਾਉਣਾ}} once, checking the whole word
\end{itemize}

\begin{tcolorbox}[breakable,colback=teal!4,colframe=teal!35!black,title={Before you move on}]
What does \pa{ਪਛਤਾਉਣਾ} mean? (\textquotedblleft{}To regret\textquotedblright{}.) Say it once more.
\end{tcolorbox}
\section[\texorpdfstring{\pa{ਗੁਜ਼ਾਰਨਾ} (guzārnā)}{ਗੁਜ਼ਾਰਨਾ (guzārnā)}]{\pa{ਗੁਜ਼ਾਰਨਾ} (guzārnā) — to spend time}

\label{lesson:PA-C317-guzarna}



\begin{quote}
Before the new one: say the Punjabi for an era, then the Punjabi for to regret.
\end{quote}

\subsection*{You'll want to know: \pa{ਗੁਜ਼ਾਰਨਾ}}
\textbf{\pa{ਗੁਜ਼ਾਰਨਾ}} — \emph{guzārnā} — \textquotedblleft{}to spend time\textquotedblright{}.

\begin{quote}
\pa{ਗਰਮੀ} \pa{ਵਿੱਚ} \pa{ਮੈਂ} \pa{ਇੱਕ} \pa{ਮਹੀਨਾ} \pa{ਨਾਨੀ} \pa{ਦੇ} \pa{ਘਰ} \pa{ਗੁਜ਼ਾਰਿਆ।}
\end{quote}

\begin{quote}
\emph{garmī vich maiṁ ikk mahīnā nānī de ghar guzāriā.}
\end{quote}

\begin{quote}
In the summer I spent a month at my grandmother's house.
\end{quote}

Past (I …): \pa{ਮੈਂ} \pa{ਛੁੱਟੀ} \pa{ਘਰ} \pa{ਵਿੱਚ} \pa{ਗੁਜ਼ਾਰੀ}, \emph{maiṁ chhuṭṭī ghar vich guzārī}.

\begin{cousinweb}
\textbf{\pa{ਗੁਜ਼ਾਰਨਾ}} (\emph{guzārnā}) is a loan, usually traced to Persian \emph{guzār}, \textquotedblleft{}passing\textquotedblright{}, with a Punjabi verb ending.
\end{cousinweb}

\begin{grammarlens}[title={What you've built}]
One more word for saying what happened.
\end{grammarlens}

\subsection*{Your turn}
\begin{itemize}
\raggedright
  \item \emph{Say it:} \emph{guzārnā}
  \item \emph{Say it:} \emph{guzārnā}, once more
  \item \emph{Say it:} \emph{garmī vich maiṁ ikk mahīnā nānī de ghar guzāriā.} — In the summer I spent a month at my grandmother's house.
  \item \emph{Recall:} say the Punjabi for an era, then the Punjabi for to regret, then say \emph{guzārnā} again
  \item \emph{Write it:} copy \textbf{\pa{ਗੁਜ਼ਾਰਨਾ}} once, checking the whole word
\end{itemize}

\begin{tcolorbox}[breakable,colback=teal!4,colframe=teal!35!black,title={Before you move on}]
What does \pa{ਗੁਜ਼ਾਰਨਾ} mean? (\textquotedblleft{}To spend time\textquotedblright{}.) Say it once more.
\end{tcolorbox}
